\chapter*{Objective}
\addcontentsline{toc}{chapter}{Objective}

The objective of Phase 2 is to convert the completed Phase 1 direction-of-arrival (DOA) sector-classification prototype from a controlled feasibility demonstration into a rigorous robustness and comparative-evaluation study for autonomous edge sensing. Phase 1 established the full reference pipeline: an eight-element uniform linear array (ULA), sample-covariance estimation, compact beamspace feature extraction, nearest-centroid continuous-valued classification, rate-based spike encoding, and leaky-integrate-and-fire (LIF)-inspired temporal evidence accumulation. Phase 2 will preserve this architecture, its sector definitions, and its principal evaluation logic so that all new results remain directly traceable to the validated Phase 1 baseline.

The proposed work will quantify how low signal-to-noise ratio, angular proximity to sector boundaries, and controlled array-model mismatch affect three inference paths: a classical array-processing reference, the Phase 1 continuous nearest-centroid classifier, and the Phase 1 spike-based LIF-inspired classifier. Performance will be evaluated using classification accuracy, accuracy retention relative to the continuous baseline, stable-decision latency, spike activity, and robustness degradation under progressively more difficult sensing conditions.

The engineering objective is not to claim that spike-based inference replaces high-resolution DOA estimation or automatically improves raw accuracy. It is to determine, under reproducible stress conditions, whether sparse temporal inference retains sufficient directional information while providing explicit latency control and a hardware-oriented state-update structure. The final Phase 2 outcome will be a defensible operating point and experimental evidence suitable for guiding a later PYNQ, FPGA, or Vitis HLS implementation stage.